\runninghead{No Want to Them That Fear Him}
\properlane{no-want-just-in-gods-hand}{No Want to Them That Fear Him, and the Just in God's Hand}

The Gradual of this Mass promises that nothing is lacking to those who fear
the Lord, and its Offertory says that the souls of the just are in the hand of
God and only seemed to die. Both statements meet an obvious objection: the
saints the feast honours went hungry, were persecuted and were put to death.
St Basil the Great says that the good the fearer cannot lack is God himself;
St Augustine, that the goods are spiritual; St Thomas Aquinas and St Robert
Bellarmine add that God helps those who seek him in bodily need. Augustine
also glosses the Offertory's first words: the just in God's hand dwell under
his help. If the good the fearer cannot lack is God himself, the just in God's
hand have it.

\subsection{A promise that seems false}
\label{sec:no-want-question}

The Gradual sings \latin{Timete Dominum, omnes Sancti eius: quoniam nihil
deest timentibus eum}, and then, \latin{Inquirentes autem Dominum, non
deficient omni bono}. St Augustine states the difficulty as an objector would.
In the psalm the promise follows the words \latin{Divites eguerunt et
esurierunt}, ``the rich have wanted and have suffered hunger,'' and
\latin{Si ad litteram accipias, videtur te fallere}: ``if you take it according
to the letter, it seems to deceive you.'' The objector has seen many wicked rich
men die in their riches, and says of himself, \latin{Quotidie ad basilicam
surgo, quotidie genuflecto, quotidie inquiro Dominum, et nihil boni habeo},
``Daily I rise for the church, daily I kneel, daily I seek the Lord, and I have
nothing good'' (\work{Enarr.\ in Ps.}\ 33, s.~2, n.~14). The Mass sings the
promise without the half-verse about the rich, but the singers still know
saints who were poor.

The Offertory takes its verses from the chapter in which Wisdom answers the
wicked who said, ``Let us condemn him to a most shameful death'' (Wis 2:20).
The antiphon says that the just \latin{visi sunt oculis insipientium mori},
seemed to die in the eyes of the unwise, and that \latin{non tanget illos
tormentum malitiae}, the torment of malice shall not touch them. The just
whose death the wicked brought about are said to be untouched. Two questions
follow: what good the fearer is promised, if it is not bread; and what it means
to be in God's hand, if it is not to be kept from death.

\subsection{St Basil: the fear that lacks nothing}
\label{sec:no-want-basil}

St Basil, preaching on Psalm~33 (34), begins from what fear does. Unless fear
has ordered our life, holiness cannot be gained in the body, and he quotes
\latin{Confige timore tuo carnes meas}, ``Pierce thou my flesh with thy fear''
(Ps 118:120; Hebrew 119:120). Then: \latin{Timenti itaque non est inopia; hoc est,
qui per timorem ab omni indecora et absurda actione arcetur, ei nihil deest ad
quamlibet virtutem, sed perfectus est}, ``to him who fears there is no want;
that is, he who is kept by fear from every unseemly and unreasonable act lacks
nothing for any virtue, but is perfect'' (\work{Hom.\ in Ps.}\ 33, n.~6).

On the next verse he offers two accounts of the rich who lacked. The words may
teach contempt of bodily riches; or the rich are Israel, to whom belonged the
adoption and the promises (Rom 9:4--5), and who, having killed the Bread of
life, suffer the famine of hearing the word. Those from the nations who seek
the Lord \latin{non minuentur omni bono. Perfectum omnino bonum Deus ipse est,
quo non minuentur quotquot exquirunt ipsum}: ``shall not be lessened in any
good. The wholly perfect good is God himself, of whom none who seek him shall
be lessened.'' He then puts the objection as a dilemma. Whoever takes bodily
goods for the good the psalm promises must say either that the apostles did
not seek the Lord, since they lacked bodily goods, or that Scripture is false.
Neither can be said: \latin{Paulus utique et inquirebat Dominum, et nihil ei
defuit bonorum}, ``Paul certainly sought the Lord, and nothing of the goods
was lacking to him,'' though he was scourged, stoned and shipwrecked (n.~7).
On Basil's account the promise is kept in the one who fears, because what he
cannot lack is God.

\subsection{St Augustine: the spiritual goods}
\label{sec:no-want-augustine}

St Augustine answers his own objection by the kind of good the psalm means:
the goods are spiritual, \latin{in bonis spiritualibus} (n.~15). The rich who
lack, he says, lack the bread that says \latin{Ego sum panis vivus}, ``I am the
living bread,'' and he adds the beatitude of those who hunger and thirst after
justice (Mt 5:6). That gloss is his comment on the half-verse about the rich,
which the Gradual does not sing; what he says of the seekers' good is his
comment on the half-verse it does sing. Of peace he says, on a later verse of
the psalm that the Mass does not use, that it is had \latin{ex parte hic \dots\
ex toto illic}, ``in part here, wholly there'' (n.~19, on v.~15).

\subsection{St Thomas and St Robert Bellarmine: help in need}
\label{sec:no-want-thomas-bellarmine}

St Thomas Aquinas begins where Basil does: \latin{Timete \dots\ sancti, quia
nullus potest esse sanctus nisi sit timens}, ``Fear, you saints, because no
one can be a saint unless he fears,'' and those who remain in holiness still
need fear (\work{Super Ps.}\ 33, n.~10). He expounds \latin{non est inopia}
three ways: of spiritual goods, since the fear of the Lord is the treasure of
Isa 33:6; of bodily want, since the man content with what he has is not needy
(Phil 4:11--12); and \latin{Item Deus quaerentibus se subvenit in necessitate},
``God also helps those who seek him in their need.'' He meets the objection of
St Paul, ``Even unto this hour we both hunger and thirst'' (1~Cor 4:11), with
the physician who withholds food from the sick for the sake of their health:
God sometimes sends want and sometimes gives riches, as it serves salvation.

St Robert Bellarmine reads the promise in two senses. In the higher, to feel
no want ``is the essence of perfection,'' and those who seek the Lord shall
lack no good ``because, as they cling to the Supreme Good, they possess all
that is good.'' In a humbler sense, the same words teach that God will supply
his servants with temporal things and not desert them in need
(\work{Explanatio in Psalmos}, on Ps 33, his verses 9--10). He says nothing of the
vision of God.

Basil and Augustine put the promised good in God and the goods of the spirit;
Aquinas and Bellarmine add providence in bodily need. The four agree that the
fearer lacks nothing that matters. They disagree whether bodily need is
included, and the disagreement stands: Basil excludes bodily goods by the
apostles' example, and Aquinas meets the same example by the physician.

\subsection{The just in the hand of God}
\label{sec:no-want-offertory}

St Augustine glosses the Offertory's first clause in a sentence:
\latin{Justorum animae in manu Dei sunt, hoc est, habitant in adjutorio
Altissimi, et in protectione Dei coeli commorantur}, ``the souls of the just are
in the hand of God, that is, they dwell in the help of the Most High, and abide
in the protection of the God of heaven'' (\work{Enarr.\ in Ps.}\ 87, n.~5,
citing Ps 90:1; Hebrew 91:1). He says it of the just among the dead before their
resurrection, and adds that those cast out of God's hand are the ones who
believed Christ had been cast out of it.

St Bonaventure, in his commentary on Wisdom, glosses the same words the same
way and goes further. The souls are first those of the martyrs, after the
\work{Glossa}, \latin{quod etiam verum est de aliis iustis}, ``which is true
also of the other just''; the verse says souls, not bodies, for the body is
given into the hands of the wicked; and \latin{in manu Dei sunt, id est in
protectione, et ideo sunt securae}, ``they are in God's hand, that is, under his
protection, and therefore they are safe,'' with the same psalm verse and the Lord's
``no man shall pluck them out of my hand'' (Jn 10:28). His text reads
\latin{tormentum mortis}, and the torment that does not touch them is eternal
death, for temporal death is \latin{quasi umbra mortis \dots\ quidam transitus
ad vitam}, ``as it were a shadow of death, a passage to life.'' What he says
of the torment is said of the Vulgate's \latin{mortis}, not of the antiphon's
\latin{malitiae}, on which none of these commentators speaks. St
Alphonsus Liguori, meditating on the same verses with the antiphon's omission,
asks: ``If God holds fast in his hands the souls of the just, who can snatch
them from him?'', and says that at the hour of death God multiplies his help
where the danger is greater (\work{Preparation for Death}, Ninth
Consideration).

They differ over whose souls are meant. Blessed Ildefonso Schuster reads the
antiphon of the martyrs: ``The persecutor believes that the lives of the
martyrs are in his hands, but they are in the hands of God'' (\work{The
Sacramentary} V, p.~211). Alphonsus reads it of the death of every just soul,
whose enemy is the last temptation; Bonaventure holds both, the martyrs
specially and all the just truly. On Schuster's account the Offertory names
the martyrs among the saints of the feast; on Alphonsus's it names every
saint.

\subsection{Joy in the Lord they have}
\label{sec:no-want-elements}

The Introit's verse bids the just exult \latin{in Domino}, in the Lord, and its
commentators (see the Introit in Each Text in Its Setting) agree on what that
joy is not: for Augustine, joy in oneself; for Basil, joy in prosperity; for
the extract under Cassiodorus's name, joy \latin{in terrenis delectationibus},
in earthly delights. Aquinas puts it in four words, \latin{in Domino, non in
mundo}, ``in the Lord, not in the world.'' Basil gives the reason the just
rejoice, that they have the Lord; and on the Gradual's psalm he says that the
good the seeker cannot lack is God himself. On both psalms he names one good,
the Lord whom the just have.

The Gospel blesses those who hunger and thirst after justice, for they shall
be filled. Augustine sets that beatitude beside the half-verse of Ps 33 the
Mass leaves out, and Basil sets it beside the verse before the Gradual's,
\latin{Gustate}, ``O taste.''

The Collect asks for the abundance of mercy that is desired,
\latin{desideratam \dots\ propitiationis abundantiam}, and the Secret that the
gifts be for our salvation, \latin{nobis salutaria}. In the Lesson the angel
from the east forbids the angels of destruction to hurt the earth until the
servants of God are sealed. The Alleluia promises those who labour \latin{ego
reficiam vos}, ``I will refresh you,'' and the Communion promises the clean of
heart that they shall see God. The Postcommunion asks that the faithful be
\latin{muniri}, defended, by the saints' prayers. Of the Sunday's prayers, the
Collect asks release from the bonds of sin, the Secret calls what God has
given a gift \latin{immeritis}, to the undeserving, and the Postcommunion asks
that those who share divine things be not left \latin{humanis \dots\
periculis}, subject to human dangers. Either Preface ends with the praise of the
angels.

\subsection{Does the promise include bread?}
\label{sec:no-want-difficulty}

The strongest objection is Augustine's: taken by the letter, the verse seems
to deceive. The four Fathers meet it in different ways, and the difference is
not only of emphasis. If Basil is right, the Gradual promises the saints God
himself and the virtues, and the hunger of the apostles is no exception to it.
If Aquinas and Bellarmine are right, it also promises help in need, and the
hunger is God's medicine.

St Bernard answers from the command rather than the promise. Preaching on the
feast, he speaks of the verse as it was sung in the night office:
\latin{Cantavimus sanctis hac nocte, dicentes: Timete Dominum, omnes sancti
ejus; sed non istis}, ``we sang to the saints this night, Fear the Lord, all
ye his saints; but not to these,'' that is, not to the saints in glory, who no
longer fear. To us it is rightly said, for \latin{Tota interim beatitudo nostra
est timere Deum}, ``meanwhile our whole beatitude is to fear God,'' while the
saints' festivity is in exultation and praise (sermo~5, n.~3).

\begin{fourSenses}
\item[Literal.] The psalmist, whom the title sets before Achimelech, promises
those who fear the Lord no want and those who seek him no lack of any good;
Wisdom's just seemed to the unwise to die, and they are in peace.
\item[Allegorical.] The rich who lacked are Israel, which killed the Bread of
life and suffered a famine of the word, and those who seek the Lord are the
nations, who lack no good (Basil).
\item[Moral.] Fear trains the life and pierces the flesh, and the fearer lacks
nothing for any virtue (Basil); no one is holy without fear (Aquinas); our
beatitude meanwhile is to fear God (Bernard); the just rejoice in the Lord and
not in themselves (Augustine).
\item[Anagogical.] God himself is the perfect good (Basil); peace is had wholly
there (Augustine); the souls of the just dwell in God's help and protection,
and are safe (Augustine, Bonaventure).
\end{fourSenses}
